\subsection*{Критерий информативности}
Как мы уже и обсуждали, мы хотим каким"=то образом отбирать наилучший предикат. Тут немного сложнее. Надо разобраться, \textit{а что означает лучший?} По какому критерию мы так считаем? Для этого введём какую-то функцию, которая будет описывать, насколько текущая вершина информативна (насколько хорошо ей удалось разбить признаки. Что такое <<хорошо разбить признаки>> "--- это пока тоже вопрос, но в контексте решения станет понятнее):

$H$ "--- \texttt{impurity}, критерий информативности.

Критерий применяется в случае классификации и принимает большие значения, если вершина содержит большое разнообразие классов, и низкое значение, если разнообразие классов маленькое.

Что может быть критерием? Ну например ранее введённая энтропия. Помимо неё существует индекс Джини.

Как и в случае с энтропией предполагаем, что прогнозом дерева является вектор вероятностей, размерность которого совпадает с количеством классов. $i$"=й элемент вектора соответствует вероятности того, что объекты, попавшие в данный листовой узел, принадлежат $i$"=му классу. Тогда:

\[
\sum_{k=1}^K p_k (1 - p_k)
\]

В целом, в большинстве случаев, энтропия и индекс Джини выдают примерно одинаковые результаты.

\subsection*{Случай регрессии}
В случае с регрессией, вместо какого"=то там критерия информативности применяются уже знакомые $MSE$, $MAE$ и т.д. Однако стратегия немного иная.

Видите ли, у каждой функции потерь есть такое понятие как \textit{оптимальный прогноз} "--- какое значение можно выдавать, чтобы конкретная функция потерь была как можно меньше. Например, в случае с $MSE$, для минимизации ошибки достаточно выдавать среднее значение целевой переменной:
\[
\sum_{i=1}^\ell (y_i - \overline{y})^2
\]

Соответственно, когда мы объявляем вершину листовой (а учитывая гибкость настройки, мы можем добиться такого в любой момент), для объектов, которые оказались в данной вершине:
\begin{itemize}
    \item Считается среднее значение их целевых переменных;
    \item Подставляется в $MSE$.
\end{itemize}

Конечно, в точности такой алгоритм подходит только для $MSE$. Для $MAE$, например, оптимальным прогнозом является \textit{медиана}.

\subsection*{Так как выбирать предикаты?}
Мы немного отошли от темы, но это было необходимо для вывода следующей формулы. Итак, формула того самого злополучного $Q$:
\[
Q(R_m, j, t) = H(R_m) - \frac{|R_l|}{|R_m|}H(R_l) - \frac{|R_r|}{|R_m|}H(R_r) \rightarrow \max
\]

т.е. мы сравниваем информативность текущей вершины с усреднёнными информативностями дочерних вершин. Если по сравнению с исходной информативностью, в какой-то из дочерних вершин получилось лучше "--- то мы полагаем, что предикат подобран хорошо. Поэтому уместна максимизация.